\documentclass[11pt,a4paper]{article}

\usepackage[margin=1in]{geometry}
\usepackage[T1]{fontenc}
\usepackage[utf8]{inputenc}
\usepackage{hyperref}
\usepackage{longtable}

\title{BigDataMind Project Progress Log}
\author{Project journal}
\date{Last updated: 8 November 2026}

\begin{document}
\maketitle

\section*{Purpose}
This document is the running record of project progress. Each time project work
is authorized and performed, add a dated entry describing the changes made,
the findings, and how the work was verified. Do not treat this log as
authorization to begin a new implementation milestone.

\section{Project overview}
BigDataMind is a planned analytics assistant for public NYC TLC yellow taxi
trip data. The intended pipeline uses Spark and a Medallion architecture:
unchanged raw inputs, Bronze ingestion, cleaned Silver data, and aggregated Gold
tables. Future milestones cover orchestration, a guarded LLM question-answering
layer, a backend, a frontend, and evaluation. Spark is responsible for data
calculations; an LLM is not to invent or calculate result values.

\section{Progress entries}

\subsection{8 November 2026 --- Environment and initial project setup}
\textbf{Reported status:} The developer completed WSL2 Ubuntu, Git, and virtual
environment configuration. A project data directory is linked to
\texttt{/mnt/d/BigDataMind-data}, keeping the large raw inputs on the D: drive.
This setup was reported by the developer and was not reconfigured during the
work recorded here.

\textbf{Changes:} No project code changes are recorded for this setup stage.

\textbf{Findings:} The repository is organized for the planned Spark, LLM,
backend, frontend, notebooks, tests, and documentation work. At the start of
the recorded project implementation, the tracked Spark implementation
consisted of a small Spark smoke-test script; the planned Bronze ingestion
module was not yet present.

\subsection{8 November 2026 --- January source data acquired}
\textbf{Changes:} The developer downloaded the January 2024 yellow taxi Parquet
file, the taxi zone lookup CSV, and the TLC yellow trip-record data dictionary
PDF into \texttt{data/raw/}. The data directory is ignored by Git.

\textbf{Findings:} The January Parquet file is approximately 48 MiB and contains
2,964,624 rows across 19 columns. The zone lookup has columns
\texttt{LocationID}, \texttt{Borough}, \texttt{Zone}, and
\texttt{service\_zone}. The source Parquet and lookup file were retained
unchanged.

\subsection{8 November 2026 --- January data exploration (Milestone 1.3)}
\textbf{Changes:} Added the read-only PySpark profile
\texttt{notebooks/01\_data\_exploration.py}. Recorded its findings in
\texttt{docs/journal.md} and provisional data-handling decisions in
\texttt{docs/decisions.md}; updated \texttt{ROADMAP.md}. These changes were
committed as \texttt{94e19f2}
(\texttt{feat: add January data exploration milestone}).

\textbf{Findings:}
\begin{itemize}
  \item The Parquet schema contains pickup and drop-off timestamps, location and
        vendor identifiers, passenger and payment fields, and trip/fare measures.
  \item There are 18 pickup timestamps outside January 2024: 15 before
        January 1 (including dates in 2002 and 2009), and 3 on February 1.
  \item There are 140,162 rows with \texttt{payment\_type=0}. Those same rows
        contain nulls in \texttt{passenger\_count}, \texttt{RatecodeID},
        \texttt{store\_and\_fwd\_flag}, \texttt{congestion\_surcharge}, and
        \texttt{Airport\_fee}. No other source columns had nulls.
  \item There are 870 nonpositive-duration trips, 60,371 zero-distance trips,
        and no negative distances.
  \item There are 37,448 negative fares and 35,504 negative total amounts.
  \item The median and 99th-percentile trip distances are 1.68 and 20.0. The
        maximum is 312,722.3; 23 trips exceed 1,000 distance units.
  \item There are no exact duplicate full rows. Each of the 265 distinct zone
        lookup IDs matches the trip data's pickup and drop-off location IDs.
\end{itemize}

\textbf{Interpretation:} These are profiling signals, not approved cleaning
rules. In particular, negative amounts may need interpretation as adjustments
or refunds; nulls are not to be imputed without evidence; unusual distances and
boundary dates are not to be removed using arbitrary thresholds. The TLC data
dictionary and reporting scope still need to inform those decisions.

\textbf{Verification:} The saved exploration script compiled and ran with the
project virtual environment and local PySpark 3.5.3. It reproduced the row
count, quality signals, zero exact duplicates, and zero unmatched zone IDs.
The profile only reads from \texttt{data/raw/} and does not write to it.

\section{Current status and next authorized step}
Milestones 1.1 (environment, as reported), 1.2 (download one month of source
data), and 1.3 (explore the data) were recorded as complete at the preceding
update. At that time, milestone 2.1 had not started. The later Bronze progress
entry below supersedes that status: milestone 2.1 is now complete. Silver
milestone 2.2 is the proposed next milestone and has not started. Continue
implementation only after the developer explicitly asks to proceed.

\section{Update convention}
Append a dated subsection after each authorized work session. Include:
\begin{enumerate}
  \item the milestone or task;
  \item files or behavior changed;
  \item observations, results, or issues found;
  \item verification performed and its outcome; and
  \item any decisions deferred to the developer.
\end{enumerate}
Do not silently rewrite earlier findings; add a dated correction if new
evidence changes them.

\section{8 November 2026 --- Bronze ingestion (Milestone 2.1)}
\textbf{Changes:} Added \texttt{spark/config.py} for data paths based on
\texttt{BDM\_DATA\_ROOT}, \texttt{spark/session.py} for the shared local Spark
configuration, and \texttt{spark/ingestion.py} for monthly raw-to-Bronze
Parquet ingestion. Added a session-scoped Spark fixture and five ingestion
tests; declared pytest in \texttt{requirements.txt}. Updated the roadmap,
journal, and decisions log. The raw input files and the existing progress
entries were not rewritten.

\textbf{Findings:} The January source schema matched the exploration baseline.
The source and Bronze row counts were both 2,964,624. The month partition
\texttt{source\_month=2024-01} contains three Parquet part files, and all 18
pickup timestamps outside January remain in Bronze. The original 19 columns
are preserved, followed by only \texttt{ingestion\_ts},
\texttt{source\_file}, and \texttt{source\_month}.

\textbf{Verification:} \texttt{.venv/bin/python -m pytest -q} completed with
\texttt{5 passed}. The command
\texttt{.venv/bin/python -m spark.ingestion --year 2024 --month 1} was run twice;
both runs logged source and Bronze counts of 2,964,624 and passed row-count
verification. A separate Spark check confirmed that source and Bronze each
contain 18 out-of-month pickup timestamps. The raw Parquet and zone CSV sizes
and modification times match the values recorded before ingestion.

\textbf{Decisions and deferred work:} Bronze derives its partition month from
the source file, partitions by \texttt{source\_month}, preserves source names
and types, and uses dynamic partition overwrite for safe reruns. Schema
differences are logged rather than treated as automatic failures. No cleaning
rules were applied; decisions about negative amounts, payment type 0,
out-of-range dates, and extreme distances remain for Silver. Silver and later
milestones have not been started.

\end{document}
